(1) 自旋半粒子初态为 $|+z\rangle$，沿 $x$ 轴测量 $\sigma_x$。求测得 $+1$ 与 $-1$ 的概率及对应测后态。
(2) 如果第一测量结果为 $+x$，随即测量 $\sigma_z$，求 $+z$ 的条件概率。
(3) 忽略第一测量读数而只保留非选择性测量作用，再测 $\sigma_z$，求最终 $+z$ 概率。
(4) 将非选择性第一测量换成无任何中间测量，比较最终结果并说明为何不相容测量会破坏原初确定性。
